La movilidad urbana constituye uno de los pilares fundamentales del desarrollo social, económico y territorial de las ciudades contemporáneas. En los entornos urbanos actuales, caracterizados por un crecimiento acelerado de la población, la expansión de las redes de transporte y el aumento constante del parque automotor, la circulación de personas y vehículos se ha convertido en un proceso cada vez más complejo. Esta complejidad no solo está relacionada con la infraestructura vial o con los sistemas de transporte disponibles, sino también con la interacción permanente entre diversos actores viales, tales como conductores, motociclistas, ciclistas, peatones y sistemas de transporte público, quienes comparten el mismo espacio urbano bajo dinámicas de movilidad intensas y cambiantes \cite{mina2024}.\\
En este escenario, la seguridad vial emerge como un desafío central para la gestión urbana y la planificación del tránsito. Los accidentes de tránsito representan una problemática de carácter multidimensional que involucra factores humanos, técnicos, ambientales y organizacionales. La ocurrencia de siniestros viales no solo implica daños materiales o afectaciones momentáneas en la movilidad, sino que genera impactos profundos en términos sociales, económicos y sanitarios. Cada accidente puede traducirse en pérdidas humanas irreparables, lesiones incapacitantes, afectaciones familiares y costos significativos para los sistemas de salud y para la economía de las ciudades. Por esta razón, la accidentalidad vial ha sido reconocida a nivel internacional como uno de los principales retos de seguridad pública en contextos urbanos \cite{valencia2025}.\\
Las grandes ciudades enfrentan particularmente este desafío debido a la alta concentración de actividades económicas, la diversidad de medios de transporte y la intensidad de los flujos vehiculares que circulan diariamente por su infraestructura vial. Medellín, como una de las principales áreas metropolitanas de Colombia, refleja con claridad estas dinámicas urbanas. La ciudad se caracteriza por una movilidad compleja influenciada por su topografía, por la alta presencia de motocicletas y por la coexistencia de diferentes sistemas de transporte público y privado. Estas condiciones generan escenarios de interacción vial constantes que incrementan la probabilidad de ocurrencia de accidentes, especialmente en corredores estratégicos, intersecciones o zonas con alta densidad vehicular \cite{ramirez2024}.\\
En este contexto, comprender el fenómeno de la accidentalidad vial exige ir más allá de la simple recopilación de registros de accidentes. En Medellín, el conjunto de datos abiertos Mede concentra registros administrativos de incidentes y víctimas con detalle temporal, territorial y por tipo de evento; no obstante, la consulta habitual de esos archivos no garantiza por sí sola una lectura integrada del territorio ni proyecciones comparables bajo criterios explícitos \cite{atancuri2024,castrillon2025}. La falta de integración entre visualización geoespacial, indicadores descriptivos y modelos orientativos sigue limitando la capacidad de las instituciones y de la ciudadanía informada para interpretar patrones de riesgo con trazabilidad en los cálculos.\\
La transformación digital de las ciudades plantea nuevas oportunidades para abordar esta problemática mediante el uso de tecnologías de información y análisis de datos. En el contexto de las llamadas ciudades inteligentes, el uso de plataformas digitales, sistemas de información geoespacial y herramientas de visualización interactiva permite convertir grandes volúmenes de datos en conocimiento útil para la gestión pública. Estas herramientas facilitan la identificación de patrones espaciales y temporales en fenómenos urbanos complejos, permitiendo a las autoridades y a los investigadores comprender con mayor precisión cómo, dónde y cuándo ocurren determinados eventos dentro del territorio \cite{changotasig2023}.\\
Particularmente, los sistemas de análisis geoespacial han demostrado ser instrumentos altamente efectivos para el estudio de fenómenos relacionados con la movilidad urbana. La georreferenciación de eventos, combinada con herramientas de visualización como mapas de calor, gráficos dinámicos y tableros de monitoreo, permite identificar zonas críticas de accidentalidad, analizar tendencias históricas y comparar el comportamiento de diferentes variables asociadas a los siniestros viales. Este tipo de análisis no solo facilita la interpretación de la información, sino que también contribuye a orientar estrategias de prevención más precisas y focalizadas \cite{monar2025}.\\
En línea con el objetivo del trabajo de grado ---\textit{Sistema de mitigación de accidentes de tránsito para ciudad con grandes urbes: caso Medellín}---, este documento presenta el diseño, la implementación y la evaluación de \textbf{ViaData — Medellín}, plataforma web desarrollada para organizar y explorar la accidentalidad registrada en Mede (periodo aproximado 2014--2021). El sistema articula ingeniería de software, análisis de datos urbanos y sistemas de información geográfica en un mismo entorno consultable.

\clearpage

\subsection{Propuesta y alcance de ViaData — Medellín}

ViaData — Medellín es una aplicación web compuesta por un backend en Django con API REST, una base PostgreSQL ampliada con PostGIS y un frontend en React. Los datos se incorporan mediante un flujo ETL reproducible que depura el Excel fuente, normaliza catálogos y carga incidentes y víctimas con coordenadas y atributos territoriales. Sobre esa base, la plataforma ofrece cuatro frentes de consulta diferenciados por rol.

El \textbf{tablero de indicadores} y el \textbf{mapa analítico} están disponibles de forma pública. El primero resume evolución mensual, distribución por día y hora, rankings territoriales y variables de gravedad; el segundo muestra puntos de incidentes, capas de densidad, coroplética por comuna o barrio y concentraciones espaciales calculadas en el servidor. Para perfiles \textit{analista} y \textit{administrador}, el módulo de \textbf{predicciones} compara modelos estadísticos transparentes ---regresión lineal, estacionalidad, Poisson, media móvil, bandas $\mu \pm 3\sigma$, ARIMA y SARIMA--- con validación \textit{hold-out}, e integra priorización territorial, carga esperada por territorio, proporción de víctimas fatales y patrones de concentración temporal. Complementan la oferta los \textbf{reportes imprimibles} y un \textbf{asistente conversacional} basado en Google Gemini que invoca las mismas rutas de la API del sistema en lugar de inventar cifras; el rol administrador además gestiona usuarios y permisos.

Las proyecciones se entienden como escenarios orientativos bajo estabilidad del patrón histórico filtrado: apoyan la planificación y la lectura comparativa del territorio, pero no establecen causalidad ni estiman probabilidad individual de sufrir un siniestro. Del mismo modo, el alcance de esta versión no incluye ranking automatizado por vía o punto crítico ---porque Mede no alimenta esos catálogos---, ni modelos de aprendizaje automático espacial avanzado; la ingesta de datos nuevos permanece documentada como proceso manual reproducible. Estas decisiones de cronograma se detallan en las secciones de metodología y trabajo futuro.

En conjunto, ViaData — Medellín busca convertir registros administrativos dispersos en un instrumento geoespacial actualizable que facilite identificar zonas críticas, seguir tendencias y discutir escenarios de carga futura con criterios explícitos, en sintonía con enfoques recientes que vinculan SIG, tableros dinámicos y gobernanza del dato en seguridad vial \cite{carrillo2024,lugo2025}.




\section{ANTECEDENTES}
En los estudios recientes sobre movilidad urbana y seguridad vial se evidencia que los sistemas geoespaciales reducen tiempos de análisis y mejoran la focalización de intervenciones. Investigaciones de los últimos años han mostrado que la identificación de zonas calientes mediante técnicas de densidad espacial y agrupación territorial permite priorizar corredores críticos antes de invertir en medidas generales \cite{xu2022,li2023}.\\\\
El trabajo desarrollado por Viveros (2022) analiza de manera rigurosa los factores que influyen en los niveles de accidentalidad en zonas de convergencia y divergencia de glorietas en la ciudad de Medellín. La investigación parte del reconocimiento de que las intersecciones, especialmente las glorietas, concentran altos índices de siniestralidad debido a la interacción simultánea de múltiples flujos vehiculares y a la incidencia del factor humano. Desde un enfoque cuantitativo, el autor recopila y clasifica datos históricos de accidentes según severidad y frecuencia, identificando intervalos críticos de tiempo con mayor ocurrencia, empleando modelos microscópicos calibrados en el software PTV-Vissim, integrando variables geométricas, volúmenes de tránsito y dispositivos de control. La calibración se realiza mediante el estadístico GEH, garantizando confiabilidad en la simulación. Los resultados evidencian que variables como el tráfico promedio anual, el número de carriles y la longitud de las rampas inciden significativamente en las tasas de accidentes. El estudio concluye que la modelación vial constituye una herramienta técnica sólida para la toma de decisiones en planificación urbana y seguridad vial, aportando evidencia cuantificable para la mitigación de riesgos en puntos críticos de la ciudad \cite{viveros2022}.\\\\
Un enfoque aplicado en Medellín reconoce patrones de accidentalidad a partir de datos espaciales y su representación mediante mapas de calor. La propuesta subraya que los siniestros no se distribuyen al azar, sino que se concentran en corredores e intersecciones donde confluyen infraestructura, comportamiento y flujo vehicular. El uso de herramientas geoespaciales permite identificar zonas críticas, comparar periodos y observar persistencia de puntos calientes, lo que orienta controles, señalización o rediseño vial en lugar de medidas uniformes. En términos prácticos, este antecedente respalda que una plataforma con capas dinámicas y filtros territoriales ---como la que se implementa en ViaData — Medellín--- puede servir de puente entre el registro administrativo y la decisión pública \cite{monar2025}.\\\\
La tesis realizada por Cristian Orihuela, aborda la siniestralidad vial desde una lógica de relación entre riesgos (factor humano, componente vehicular y condición vial) y el comportamiento del fenómeno en un territorio urbano específico. El estudio, cuantitativo y de corte transversal, utiliza encuesta como instrumento principal y reporta una relación fuerte entre factores de riesgo y ocurrencia de accidentes, lo que refuerza la idea de que la accidentalidad puede explicarse y gestionarse si se identifican variables clave y se monitorean de forma sistemática. Más allá del resultado estadístico, el documento aporta un hilo argumental útil para proyectos de ingeniería de software: cuando los riesgos se miden de manera continua y comparable, se puede orientar la intervención pública con criterios verificables. En esa línea, una plataforma en tiempo real sería coherente con la necesidad de pasar de reportes tardíos a analítica oportuna que permita prevenir, focalizar recursos y evaluar impactos con indicadores \cite{orihuela2025}.\\\\
Un estudio, desarrollado en Bogotá por Ramirez \& Caicedo, (2024), analiza la siniestralidad en un grupo laboral con alta exposición vial (motociclistas) y describe cómo el problema no se limita a un solo factor: se combinan comportamientos (infracciones, velocidad), condiciones del entorno (vías) y elementos organizacionales. Metodológicamente integra registros e instrumentos de percepción, lo que evidencia que el dato de accidentes suele ser insuficiente si no se complementa con información contextual para comprender causas repetidas. En clave de software, este antecedente sugiere que el valor no está solo en almacenar reportes, sino en estructurar variables (tipo de actor vial, conducta, estado de vía, lugar) y habilitar análisis comparativos para intervenir sobre lo prevenible. Así, una plataforma con analítica espacial y tableros puede contribuir a priorizar acciones, enfocando campañas, controles y mejoras de infraestructura donde el riesgo se repite \cite{ramirez2024}.\\\\
La investigación evidencia que la calidad del dato es una condición previa para cualquier política efectiva de seguridad vial. La revisión sistemática concluye que, sin un sistema de indicadores estructurado y sostenido, resulta difícil estimar riesgos, priorizar interventions y evaluar si las medidas realmente salvan vidas. El valor del trabajo está en su énfasis sobre variables que deben registrarse y estandarizarse: víctimas (mortalidad/morbilidad), temporalidad, información geoespacial, características sociodemográficas y causas/tipos de siniestro. Para un trabajo de grado en ingeniería de datos y software, este antecedente es directamente transferible: una plataforma en Medellín debe diseñarse desde el modelo de datos y la gobernanza de la información, garantizando consistencia, trazabilidad y comparabilidad. Además, la revisión sugiere que la utilidad pública depende de que el sistema sea operable en el tiempo, con acceso ágil a estadísticas y con capacidad de alimentar decisiones, no solo de producir reportes aislados. Por ello, una plataforma con mapas de calor, gráficos y estadísticas cobra sentido como infraestructura de conocimiento: convierte el registro en inteligencia para prevenir y mitigar siniestros \cite{oviedo2025}.  
